\section{Data Appendix: Quantitative Synthesis and Qualification Protocols}
\label{sec:appendix}

This appendix consolidates the quantitative evidence behind Sections~\ref{sec:corrosion}--\ref{sec:applications} and specifies the test protocols a qualification program would actually run. Where the literature is heterogeneous, we report ranges and state the reasons formal meta-analysis is not yet possible.

\subsection{A1. Quantitative electrochemical synthesis vs.\ LPBF 316L}
\label{sec:a1}

Table~\ref{tab:electrochem} compiles the principal electrochemical data for the main AM-HEA families in chloride media, with 316L benchmarks. Three cautions govern interpretation. \emph{First}, reference electrodes differ between studies (SCE, Ag/AgCl, saturated calomel); values are quoted as published and can shift by tens of mV when converted, so cross-study comparisons of absolute potentials carry an uncertainty of that order. \emph{Second}, surface finish and sweep protocol are not uniform; reported $i_{\mathrm{corr}}$ for the same alloy family spans roughly an order of magnitude across studies. \emph{Third}, build quality is a hidden variable: the same nominal alloy can rank above or below cast material depending on defect content (Section~\ref{sec:a2}). Consequently, honest error bars cannot be assigned today; the ranges in Table~\ref{tab:electrochem} are the defensible summary statistic.

The key magnitudes: corrosion current densities of well-built CoCrFeMnNi-type deposits cluster around 0.09--0.13~$\mu$A/cm$^{2}$ in aerated 3.5~wt\% NaCl, with pitting potentials of order +0.1 to +0.2~V versus common references \cite{met2019slmcorr}; one SLM study reports $E_{\mathrm{pit}}=197$~mV for AM against only 37~mV for cast, with passivation windows of 386 vs.\ 216~mV \cite{met2019slmcorr}. For the LPBF 316L reference itself, AM material typically \emph{exceeds} wrought: critical pitting temperature of 27.5--31~$^{\circ}$C (build-direction dependent) against 16~$^{\circ}$C for wrought in 1~M NaCl, and pitting potentials some 330~mV nobler at 25~$^{\circ}$C \cite{slm316l_cpt}. This sets a high bar: an AM HEA must beat a strong AM-316L baseline, not just a wrought one. EHEAs passivate readily in neutral NaCl (AlCoCrFeNi$_{2.1}$ between 5 and 45~$^{\circ}$C), but lose passivity completely in the presence of reducing bisulfite---a proxy for occluded/crevice chemistry \cite{jac_ehea2023}---and Mo additions lift the pitting potential to $+282$~mV with low passive current \cite{mo_ehea_2026}. Charge-transfer resistances for passive CoCrFeMnNi-type surfaces are of order $10^{5}$~$\Omega\cdot$cm$^{2}$, comparable to 304 SS \cite{fs_coating2026}.

A minimum reporting standard for future work should follow: potentials converted to and reported against Ag/AgCl/seawater; surface preparation stated per ASTM G1 \cite{astm_g1} (e.g., 600-grit SiC finish, ultrasonic clean); sweep conventions per ASTM G3/G5 \cite{astm_g3,astm_g5}; $n\geq3$ specimens per condition with spread reported; and reference alloys (LPBF 316L, NiAl bronze) tested in the same cell.

\subsection{A2. The LPBF CoCrFeMnNi vs.\ LPBF 316L benchmark: conditions and robustness}
\label{sec:a2}

The widely quoted ``$\sim$2.2$\times$ better in chloride, $\sim$15.5$\times$ worse in acid'' result comes from a single study in the journal \emph{High Entropy Alloys \& Materials} \cite{lpbf_cantor_2024}. From the published record: specimens were printed in two orientations (perpendicular and parallel to the build direction) at an optimized volumetric energy density (window 60--110~J/mm$^{3}$, optimum 104~J/mm$^{3}$); porosity was quantified by X-ray micro-CT; both as-built and HIPed conditions were tested; media were 3.5~wt\% NaCl and 1~N H$_{2}$SO$_{4}$, with 316L and 430 SS as references; HIPed material retained the chloride advantage while hardness fell from $229\pm5$ to $152\pm4$~HV \cite{lpbf_cantor_2024}. Surface preparation for the electrochemistry is not fully documented in the open record---a gap reviewers should note; standard practice is grinding to at least 600-grit per ASTM G3/G5 conventions \cite{astm_g3,astm_g5}.

\emph{Robustness is not yet demonstrated.} No multi-lot powder or multi-build replication has been published for this benchmark, and the wider literature contains direct contradictions: one SLM-vs-cast study found the \emph{cast} alloy more resistant in 3.5~wt\% NaCl because its SLM builds carried higher densities of pores, cracks and dislocations ($i_{\mathrm{corr}}$: cast $1.20\times10^{-7}$ vs.\ XY $1.23\times10^{-7}$ vs.\ XZ $1.33\times10^{-7}$~A/cm$^{2}$; $E_{\mathrm{corr}}$ vs.\ SCE: $-0.264$, $-0.346$, $-0.323$~V respectively) \cite{frontiers2020slmcast}, while SLM CoCrNi has likewise been reported inferior to cast in one study and superior in another \cite{laam_cocrni2021}. The reconciling variable is defect content, not chemistry alone (Section~\ref{sec:a3}). Build-orientation effects on $E_{\mathrm{corr}}$/$i_{\mathrm{corr}}$ of a few tens of mV and $\sim$10--20\% are consistently observed \cite{frontiers2020slmcast,slm_params2024}. A qualification-grade version of the benchmark therefore requires: $\geq2$ powder lots, $\geq3$ builds with different orientations per condition, documented surface preparation, and both NaCl and crevice-type media.

\subsection{A3. Defect classes and pit initiation statistics}
\label{sec:a3}

Quantitative defect--corrosion links exist mainly from AM stainless steels and transfer to HEAs by analogy (Table~\ref{tab:defectpit}). The strongest anchors: (i) LPBF 420 SS with lack-of-fusion (LOF) pores averaging 18~$\mu$m at 0.42\% porosity developed $\sim$100~$\mu$m pits at defect zones versus $<$2~$\mu$m pits in defect-free zones---a roughly 50$\times$ pit-size amplification from a few tenths of a percent of porosity \cite{npd_am_review2026}; (ii) for SLM 316L, total porosity in the 0.01--0.4\% range did not shift the pitting potential but increased metastable pitting and lowered the repassivation potential, i.e., porosity governs initiation statistics and stability of propagation rather than the thermodynamic pitting threshold \cite{iop2024def304l}; (iii) minimizing LOF in DED 304L significantly improved global corrosion response \cite{iop2024def304l}; and (iv) LOF pores are structurally complex---oxide-film-lined cavities with nanocrystalline shells, predominantly spatter-derived---so each surface-breaking LOF is a pre-formed crevice with occluded chemistry \cite{lof2021}. MnS-type inclusions, the classic pit initiators in wrought stainless, are refined or absent in LPBF material, which is why even rough as-built LPBF 316L often outpits wrought \cite{jom_pitting_review2022}; in WAAM HEAs, by contrast, MnS and oxide inclusions have been observed and must be counted \cite{andersson2026waam}. Location effects (build-plate position relative to shielding-gas flow) measurably change pit initiation and should be mapped \cite{iop2024def304l}.

For defect acceptance criteria we recommend adapting fatigue-style criteria to corrosion service: (a) near-surface ($<$200~$\mu$m depth) LOF classified separately from interior porosity, since only surface-connected defects initiate external pitting; (b) maximum-pore-size limits rather than volume-fraction-only limits, given the pit-amplification evidence \cite{npd_am_review2026}; (c) inclusion density limits for WAAM deposits; and (d) statistical validation using the small-electrode method---measuring initiation-site density as a function of electrode area, as demonstrated for sintered 316L---to obtain pit-initiation \emph{statistics} rather than single-curve pass/fail outcomes \cite{npj_sintered2024}.

\subsection{A4. Galvanic protocols and cathodic protection compatibility}
\label{sec:a4}

For cladding and repair applications, galvanic behavior must be tested, not inferred. Recommended protocol per ASTM G71 \cite{astm_g71}: duplicate or triplicate HEA--substrate couples (HEA on AH36/marine steel; HEA on NiAl bronze for propeller repair) in ASTM D1141 seawater \cite{astm_d1141} and, where possible, natural seawater; area ratios spanning the service extremes (small HEA repair on large steel structure $\Rightarrow$ large cathode/small anode ratio reversed and vice versa); continuous current logging for 30--90 days plus mass-loss validation; and free-corrosion potential mapping of each material before coupling.

Acceptance thresholds: no universal numerical criterion exists for HEA couples, so we propose a two-part engineering criterion: (i) the galvanic current density shall not raise the substrate's corrosion rate by more than 10\% of its uncoupled value (or below the design corrosion allowance rate, whichever is stricter); and (ii) no localized attack shall initiate at the HEA/substrate interface or at coating defects during the test. Passive FCC HEAs are expected to be noble to structural steel, so the protected member is the substrate at defects---exactly the geometry where acceptance matters.

Interaction with cathodic protection (CP): offshore CP design typically holds carbon/low-alloy steel at or below about $-0.80$~V vs.\ Ag/AgCl/seawater with sacrificial anodes, and limits potentials more negative than roughly $-1.10$~V vs.\ Ag/AgCl to avoid hydrogen damage to higher-strength components \cite{dnv_rpb401}. The qualification questions for an HEA in a CP system are therefore: does the HEA remain passive (or behave acceptably) at $-0.80$~V imposed potential; does it inflate the CP current demand of the assembly; and, for any martensitic or high-strength HEA regions (e.g., in graded deposits), is hydrogen uptake acceptable---one LPBF CoCrFeMnNi study reports hydrogen permeation about five times lower than 316L, a favorable indicator \cite{lpbf_cantor_2024}, but imposed-potential testing in D1141 at the CP set point is required evidence.

\subsection{A5. Minimum natural-seawater exposure matrix}
\label{sec:a5}

ASTM G52 provides the governing practice for surface-seawater exposure, covering full immersion, tidal and splash zones, and spray, with crevice configurations per Guide G78 and control materials per Test Method D3623 \cite{astm_g52,astm_g78}. Table~\ref{tab:exposure} specifies a minimal matrix we judge persuasive for an initial class-society submission. Design principles: at least two exposure sites with different temperature/salinity/biota; triplicate specimens per cell; two surface conditions (as-built and machined/polished) and two build orientations for AM material; crevice assemblies on every retrieval interval; and reference alloys (LPBF and wrought 316L, NiAl bronze, 2205) in identical racks. Durations of 1, 6, 12 months are the persuasive minimum, with 24-month follow-ups for any component seeking full class approval. Endpoints per interval: mass loss per ASTM G1/G31 \cite{astm_g1,astm_g31}, maximum pit depth and rating per ASTM G46 \cite{astm_g46}, retained tensile properties, crevice inspection, and biofilm sampling (Section~\ref{sec:a9}). Note that G52 explicitly excludes high-velocity service; pumped-seawater velocity loops (0.5--3~m/s) must be added separately for heat-exchanger and pump qualification.

\subsection{A6. Integrating computational screening with printability}
\label{sec:a6}

The components of a closed-loop design pipeline now all exist. \emph{Printability:} a demonstrated high-throughput framework couples CALPHAD thermo-physical properties (solidus/liquidus, Scheil paths) with fast melt-pool surrogate models and physics-based defect criteria for lack-of-fusion, balling, keyholing and hot cracking, producing probabilistic printability maps; for the equiatomic CoCrFeMnNi space the defect-free processing fraction (printability index) was 17--22\%, and Co--Fe- and Cr--Co-rich sub-regions offered the best trade-off between balling and hot-crack resistance \cite{npm_printability2025}. Hot-crack susceptibility is quantified by Kou's crack susceptibility index (maximum of $|\mathrm{d}T/\mathrm{d}f_s^{1/2}|$ near the terminal solidification stage) computed from CALPHAD Scheil paths, an approach validated by weldability data for refractory systems and extendable to HEA databases \cite{lanl_csi2025}. \emph{Passivity:} composition-to-corrosion surrogates can be trained on the growing polarization datasets for HEAs \cite{qiu2017,qiu2019}, augmented by DFT-derived dissolution and oxide-stability descriptors and by SKPFM work-function contrasts between phases as micro-galvanic proxies \cite{lu2018design,quiambao2019}. \emph{Integration:} multi-objective Bayesian or genetic-algorithm optimization over (printability index, CSI, passivity proxy, cost/critical-element penalty) produces a Pareto front of candidate compositions; hybrid ML--CALPHAD frameworks are the current best practice for such coupled objectives \cite{jmatman_ml2026}. \emph{Validation:} compositionally graded DED libraries or LPBF combinatorial coupons test the top candidates in days, feeding back to the surrogates. WAAM screening differs mainly in feedstock constraints (wire manufacturability, evaporation loss per Section~\ref{sec:a8}) entering the objective function.

\subsection{A7. Post-processing balance for EHEAs}
\label{sec:a7}

Table~\ref{tab:postprocess} summarizes the evidence. The controlling fact for EHEAs is the heat-treatment trade-off: SLM AlCoCrFeNi$_{2.1}$ in the as-built state has the most homogeneous inter-phase chemistry and the best chloride corrosion resistance, while annealing restores FCC/B2 segregation and dramatically degrades passivity \cite{slm_ehea_2025}. HIP is the safest densification route because it closes porosity without deliberately driving segregation: HIPed LPBF CoCrFeMnNi kept chloride performance equal to as-built \cite{lpbf_cantor_2024}, at the cost of softening ($229\to152$~HV). For surface finishing of wetted faces, 316L analogues are strong: CW laser polishing (remelting mode) cut roughness by 85\% (Sa 2.1$\to$0.3~$\mu$m), reduced $i_{\mathrm{corr}}$ by 97\% and widened the passive window by 98\%, whereas nanosecond polishing introduced microcracks and femtosecond polishing gave limited benefit \cite{laserpolish2026,laserpolish2023}; electropolishing reaches internal channels, enriches the passive film in Cr (Cr/Fe ratio roughly tripled), and cut $i_{\mathrm{corr}}$ from 2.5 to 0.8~$\mu$A/cm$^{2}$ in 316L \cite{ep2014}. Recommended EHEA sequence: HIP $\rightarrow$ machining to near-final geometry $\rightarrow$ electropolish or CW laser polish of wetted surfaces $\rightarrow$ low-temperature stress relief only if required (and corrosion re-qualified). Any homogenization anneal must be justified mechanically and re-qualified electrochemically.

\subsection{A8. Mn/Cr evaporation in WAAM: magnitude and compensation}
\label{sec:a8}

Evaporation is a real but manageable effect. In cable-wire WAAM of CoCrFeMnNi with a pre-calculated chemistry, measured deviations were modest for the principal elements (Mn $\sim$2\%, Ni $\sim$10\% relative error) but large for minor, volatile, or surface-active additions (Al, Si) \cite{osintsev_waam}. Metal-cored-wire WAAM of a CoCrFeMoNiV alloy showed no clear bulk burn-off, with only slight losses of Cr, Mo and Co \cite{treutler_waam}; Cantor-type MPCW deposits exhibit minor Mn segregation and MnS/oxide inclusions \cite{andersson2026waam}. Practice guidance: (i) over-alloy Mn in the wire (welding-wire practice uses 1--3~wt\% excess) and \emph{verify deposited chemistry per build}---deposited composition is a process output, not an inference from wire chemistry; (ii) prefer low-heat-input transfer modes (CMT/short-circuit) with shorter molten-pool residence time; (iii) use trailing and side shielding or local enclosures---in reactive-alloy WAAM, supplementary shielding cut surface oxygen from $\sim$32 to $\sim$6~wt\%, illustrating the magnitude of atmosphere control available \cite{waam_mg2026}; (iv) keep Cr above the passivity threshold (order 15~at\%) after losses, since Cr supply to the passive film is the corrosion-critical variable. Corrosion-wise, mild Mn depletion is not automatically harmful---the absence of MnS inclusions is what lifts pitting resistance in LPBF 316L \cite{jom_pitting_review2022}---but it must be measured, not assumed.

\subsection{A9. Biofilm formation and under-film chemistry on HEAs}
\label{sec:a9}

Preliminary studies \emph{do} exist---this is no longer a blank field, though no AM-HEA biofilm study in natural seawater has been published yet. First-generation results: marine \emph{Pseudomonas aeruginosa} biofilms on FeCoCrNiMo$_{0.1}$ thinned and weakened the passive film and drove selective Fe dissolution via extracellular electron transfer (the first HEA MIC study) \cite{lou2020mic}; on Cantor-type FeCoNiCrMn in simulated seawater, the same organism produced maximum pit depths of 4.77~$\mu$m at 14 days against 3.16~$\mu$m sterile ($\sim$1.5$\times$), with OCP dropping from $-367$ to $\sim-500$~mV and polarization resistance falling \cite{mic_cantor2022}; and the sulfate-reducing bacterium \emph{Desulfovibrio vulgaris} caused sulfide formation, passive-film thinning and pitting on CoCrFeMnNi, becoming \emph{more} aggressive under carbon-starvation conditions that upregulate hydrogenase-mediated electron uptake \cite{wang2023srb}. The practical implication: sterile-solution screening underestimates localized attack, and occluded, nutrient-limited niches are the aggressive case.

Table~\ref{tab:biofilm} outlines a program to close the gap for AM HEAs.

\subsection{A10. Expanded standards mapping}
\label{sec:a10}

Table~\ref{tab:standards} now includes the pitting/crevice (ASTM G48, G78, G192, G46), stress-corrosion (G30, G38/G39, G129, G123), galvanic (G71), and field-exposure (G52, G31) methods proposed for the qualification ladder of Section~\ref{sec:applications}, together with the CP design reference \cite{dnv_rpb401}.

% Appendix tables
\begin{table*}[t]
\centering
\caption{Quantitative electrochemical synthesis for the main AM-HEA families in chloride media, with 316L benchmarks. Potentials are quoted as published (reference electrodes differ between studies: SCE or Ag/AgCl; absolute values shift by tens of mV under conversion). ``AB'' = as-built. Ranges across studies are the honest summary statistic; see Section~\ref{sec:a1}.}
\label{tab:electrochem}
\footnotesize
\begin{tabularx}{\textwidth}{@{}llllllX@{}}
\toprule
System & Route / state & Medium & $E_{\mathrm{corr}}$ & $i_{\mathrm{corr}}$ ($\mu$A/cm$^2$) & $E_{\mathrm{pit}}$ / key metric & Source \\
\midrule
CoCrFeMnNi & SLM, AB & 3.5\% NaCl & $-189$ mV & 0.09 & $+197$ mV; passivity window 386 mV & \cite{met2019slmcorr} \\
CoCrFeMnNi (cast ref.) & Cast & 3.5\% NaCl & $-179$ mV & 0.11 & $+37$ mV; window 216 mV; 0.91 $\mu$m/yr & \cite{met2019slmcorr} \\
CoCrFeMnNi & SLM, XY / XZ & 3.5\% NaCl & $-0.346$ / $-0.323$ V vs SCE & 0.123 / 0.133 & Cast best: $-0.264$ V, 0.120; defect-dominated build & \cite{frontiers2020slmcast} \\
CoCrFeMnNi & LPBF, AB + HIP & 3.5\% NaCl & --- & --- & $\sim$2.2$\times$ better than 316L; $\sim$15.5$\times$ worse in 1N H$_2$SO$_4$; H permeation $\sim$5$\times$ below 316L & \cite{lpbf_cantor_2024} \\
CoCrFeMnNi (friction-surfaced) & Solid-state deposit & 3.5\% NaCl & --- & 0.11 (cast: 0.55) & $+125.7$ mV (cast: $-120.3$); $R_{\mathrm{ct}}\ 7.1\times10^{5}\ \Omega\cdot$cm$^2$ & \cite{fs_coating2026} \\
CoCrNi & LAAM & 3.5\% NaCl & --- & lower with more oxides & No pitting in test window; superior to 316L & \cite{laam_cocrni2021} \\
AlCoCrFeNi$_{2.1}$ & SLM, AB & 3.5\% NaCl & --- & --- & Best of as-built/annealed/cast; annealing degrades & \cite{slm_ehea_2025} \\
AlCoCrFeNi$_{2.1}$ & Cast & 3.5\% NaCl (+NaHSO$_3$) & --- & --- & Passivates 5--45~$^{\circ}$C; \emph{no} passivation with bisulfite; pits at Al--Ni-rich B2 & \cite{jac_ehea2023} \\
AlCoCrFeNi$_{2.1}$Mo$_{0.3}$ & Cast (AM design target) & 3.5\% NaCl & --- & low passive current & $E_{\mathrm{pit}} +282$ mV; Mo/Cr composite oxide film & \cite{mo_ehea_2026} \\
316L (AM benchmark) & SLM & 1 M NaCl & --- & --- & CPT $27.5$--$31$~$^{\circ}$C vs $16$~$^{\circ}$C wrought; $E_{\mathrm{pit}} \approx +330$ mV vs wrought at 25~$^{\circ}$C & \cite{slm316l_cpt} \\
\bottomrule
\end{tabularx}
\end{table*}

\begin{table*}[t]
\centering
\caption{Defect classes and pit-initiation evidence (AM stainless steels; transfer to HEAs by analogy). Basis for proposed defect acceptance criteria in Section~\ref{sec:a3}.}
\label{tab:defectpit}
\footnotesize
\begin{tabularx}{\textwidth}{@{}lXX@{}}
\toprule
Defect class & Quantitative link to localized corrosion & Consequence for acceptance criteria \\
\midrule
Lack-of-fusion (LOF) pores & LPBF 420 SS: 18~$\mu$m avg LOF at 0.42\% porosity $\Rightarrow$ $\sim$100~$\mu$m pits at defect zones vs $<$2~$\mu$m defect-free ($\sim$50$\times$ amplification) \cite{npd_am_review2026}; LOF are oxide-lined, spatter-derived pre-formed crevices \cite{lof2021}; minimizing LOF in DED 304L significantly improves corrosion response \cite{iop2024def304l} & Class surface-near ($<$200~$\mu$m) LOF separately from interior porosity; maximum-pore-size limits, not volume fraction alone \\
Gas porosity (total volume) & SLM 316L: 0.01--0.4\% porosity does not shift $E_{\mathrm{pit}}$ but raises metastable-pitting frequency and lowers repassivation potential \cite{iop2024def304l} & Control porosity for initiation-statistics and crevice-like behavior, not just the thermodynamic pitting threshold \\
Inclusions (MnS, oxides) & MnS and LOF jointly identified as initiation sites in sintered 316L; site density quantified via electrode-area scaling \cite{npj_sintered2024}; MnS/oxides present in WAAM Cantor deposits \cite{andersson2026waam}; refined/absent MnS is why LPBF 316L can beat wrought \cite{jom_pitting_review2022} & Inclusion density limits for WAAM; exploit inclusion refinement in powder-bed routes \\
Surface roughness / as-built skin & 10--30~$\mu$m Sa typical as-built LPBF 316L; lower roughness $\Rightarrow$ higher breakdown potential; top surfaces best, downskin worst \cite{jom_pitting_review2022}; rough as-built surfaces raise SCC susceptibility \cite{lpbf_316l_scc2022} & Specify the wetted-surface condition (machined, electropolished or laser-polished) as part of the material state \\
Build-plate location & Specimens downstream of shielding-gas flow show least corrosion resistance in LPBF 304L \cite{iop2024def304l} & Map properties across the build plate during qualification, not only at coupon locations \\
\bottomrule
\end{tabularx}
\end{table*}

\begin{table*}[t]
\centering
\caption{Minimal natural-seawater exposure matrix for initial qualification of AM HEAs, structured on ASTM G52 \cite{astm_g52}, with crevice assemblies per ASTM G78 \cite{astm_g78}. All cells: triplicate specimens; two surface conditions (as-built, machined/polished); two build orientations; reference alloys (LPBF 316L, wrought 316L, NiAl bronze, 2205) in identical racks; at least two sites with different temperature/salinity/biota.}
\label{tab:exposure}
\footnotesize
\begin{tabularx}{\textwidth}{@{}lllXX@{}}
\toprule
Zone & Configuration & Durations & Flow condition & Primary endpoints \\
\midrule
Atmospheric / splash & Flat panels, 90$^{\circ}$ racks \cite{astm_g52} & 1, 6, 12, 24 mo & Wet--dry cycling, natural splash & Mass loss (G1/G31 \cite{astm_g1,astm_g31}), pit depth (G46 \cite{astm_g46}), rust/film analysis \\
Tidal & Panels on tidal rack \cite{astm_g52} & 1, 6, 12, 24 mo & Semidiurnal immersion cycles & Pit statistics, biofilm sampling, retained tensile properties \\
Full immersion (1--3 m) & Panels + crevice assemblies (G78 washers) + HEA--steel couples \cite{astm_g78,astm_g71} & 1, 6, 12, 24 mo & Quiescent to local tidal flow & Crevice attack rating, galvanic current log, EIS at retrieval intervals \\
Velocity loop (laboratory) & Pumped seawater rig & 3--6 mo & 0.5--3 m/s (pump/heat-exchanger service) & Erosion-corrosion, passive-film stability under flow \\
\bottomrule
\end{tabularx}

\vspace{4pt}
\footnotesize High-velocity service is outside G52's scope and must be added as a separate loop. Endpoints at every retrieval: mass loss, maximum pit depth and rating, crevice inspection, retained mechanical properties, and biofilm/under-film sampling (Table~\ref{tab:biofilm}).

\end{table*}

\begin{table*}[t]
\centering
\caption{Post-processing balance for EHEA and FCC-HEA marine parts: defect mitigation versus segregation/passivity trade-offs.}
\label{tab:postprocess}
\footnotesize
\begin{tabularx}{\textwidth}{@{}lXXX@{}}
\toprule
Treatment & Effect on defects/segregation & Measured effect on corrosion & Recommendation for marine EHEA/FCC parts \\
\midrule
HIP & Closes internal porosity; near-isothermal, does not deliberately drive segregation & HIPed LPBF CoCrFeMnNi keeps AB-level chloride performance; hardness $229\!\to\!152$ HV \cite{lpbf_cantor_2024} & Use when porosity is the concern; safest densification \\
Homogenization anneal & Relieves stress but restores FCC/B2 segregation in EHEAs & SLM AlCoCrFeNi$_{2.1}$: as-built best; annealing dramatically degrades chloride passivity \cite{slm_ehea_2025} & Avoid for passivity-critical parts; if required, re-qualify electrochemically \\
CW laser polishing & Localized remelting; Sa 2.1$\to$0.3~$\mu$m ($-85$\%), grain refinement $-34$\% & $i_{\mathrm{corr}}$ $-97$\%; passive window $+98$\% (LPBF 316L) \cite{laserpolish2026}; NS polishing risks microcracks, FS limited benefit \cite{laserpolish2023} & Preferred for accessible wetted faces of LPBF parts \\
Electropolishing & Removes surface peaks/reaction layer; reaches internal channels & Cr/Fe of passive film $\sim$3$\times$; $i_{\mathrm{corr}}$ 2.5$\to$0.8~$\mu$A/cm$^2$ (316L) \cite{ep2014} & Preferred for conformal/internal seawater channels \\
\bottomrule
\end{tabularx}

\vspace{4pt}
\footnotesize Recommended sequence: HIP $\rightarrow$ machining $\rightarrow$ electropolish or CW laser polish of wetted surfaces $\rightarrow$ low-temperature stress relief only if needed (with corrosion re-qualification).

\end{table*}

\begin{table*}[t]
\centering
\caption{Status of HEA biofilm/MIC evidence and a proposed program for AM HEAs. Existing studies are on cast/wrought material in simulated media; no AM-HEA natural-seawater biofilm study exists yet.}
\label{tab:biofilm}
\footnotesize
\begin{tabularx}{\textwidth}{@{}lX@{}}
\toprule
Element & Design \\
\midrule
Existing evidence & \emph{P.~aeruginosa} on FeCoCrNiMo$_{0.1}$: passive film thinned, selective Fe dissolution via EET (first HEA MIC study) \cite{lou2020mic}; on Cantor-type HEA: max pit 4.77 vs 3.16~$\mu$m (14~d), OCP $-367\!\to\!\sim\!-500$ mV \cite{mic_cantor2022}; \emph{D.~vulgaris} (SRB) on CoCrFeMnNi: sulfides, film thinning, more aggressive under carbon starvation \cite{wang2023srb} \\
Materials & LPBF CoCrFeMnNi, SLM AlCoCrFeNi$_{2.1}$, and one WAAM deposit; as-built vs machined; LPBF 316L and 2205 as controls \\
Organisms/media & \emph{P.~aeruginosa} (aerobic model), \emph{D.~vulgaris} (SRB), natural-seawater biofilm controls; 2216E and ASTM D1141 media \cite{astm_d1141}; carbon-starved variants \\
Electrochemistry & OCP, LPR and EIS time series over 14--28 days, inoculated vs sterile (per \cite{mic_cantor2022,wang2023srb} protocols); $n\geq3$ \\
Surface/film analysis & CLSM live/dead imaging, SEM pit mapping, XPS/AES depth profiles of passive film, SKPFM micro-galvanic maps \\
Under-film chemistry & pH and O$_2$ microelectrodes beneath biofilms at defined sites; sulfide assays for SRB cells \\
Integration & Couple one coupon set into the natural-seawater racks of Table~\ref{tab:exposure} so laboratory MIC and field biofilm data converge \\
\bottomrule
\end{tabularx}
\end{table*}

